% =====================================================================
%  2.3 Andregradslikninger — fasit
% =====================================================================

\fk{Del 1 — Faktorisering og nullproduktregelen}
\fo{1.1} a) $x = 1$ eller $x = -4$ \ b) $x = 0$ eller $x = 7$ \ c) $x = 3$ eller $x = -5$ \ d) $x = -2$ eller $x = 8$
\fo{1.2} a) $x = 0$ eller $x = 8$ \ b) $x = 0$ eller $x = -5$ \ c) $x = 7$ eller $x = -7$ \ d) $x = 4$ eller $x = -4$
\fo{1.3} a) $2$ og $5$ \ b) $(x+2)(x+5)$ \ c) $x = -2$ eller $x = -5$
\fo{1.4} a) Sum--produkt: $(x-5)(x+2) = 0$, så $x = 5$ eller $x = -2$ \ b) Kvadratsetning: $(x-5)^2 = 0$, så $x = 5$
c) Felles faktor: $6x(x-3) = 0$, så $x = 0$ eller $x = 3$ \ d) Flytt over, så sum--produkt: $x^2 + 2x - 15 = (x+5)(x-3) = 0$, så $x = -5$ eller $x = 3$
Andre metodevalg med god begrunnelse er riktige.
\fo{1.5} a) $x^2 - x - 20 = (x-5)(x+4) = 0$, så $x = 5$ eller $x = -4$ \
b) $x^2 + 3x - 28 = (x-4)(x+7) = 0$, så $x = 4$ eller $x = -7$
\fo{1.6} Feilen er i siste overgang. Fra $(x+5)(x-4) = 0$ får vi $x + 5 = 0$ eller $x - 4 = 0$, altså $x = -5$ eller $x = 4$.
Ola byttet fortegn. Kontroll: $x = 4$ gir $4^2 + 5 \cdot 4 = 36$ og $4 \cdot 4 + 20 = 36$.
\fo{1.7} a) Noen ganger. $k > 0$ gir to løsninger, $k = 0$ gir én ($x = 0$) og $k < 0$ gir ingen.
b) Noen ganger. Når $m \neq n$ er det to løsninger. Når $m = n$ er det bare én.
c) Alltid. Innsetting gir $0^2 + b \cdot 0 = 0$ og $(-b)^2 + b(-b) = b^2 - b^2 = 0$. For $b = 0$ er de samme tall.
\fo{1.8} $k = 12$ og $k = -12$. Kvadratsetningen gir $(x \pm 6)^2 = x^2 \pm 12x + 36$, en faktor som gjentas. $(x+6)^2 = x^2 + 12x + 36$ gir $x = -6$. $(x-6)^2 = x^2 - 12x + 36$ gir $x = 6$.
\fo{1.9} a) $(x-3)(x-8) = x^2 - 11x + 24$, så $p = -11$ og $q = 24$.
b) $(x - x_1)(x - x_2) = x^2 - (x_1 + x_2)x + x_1 x_2$. Derfor er $p = -(x_1 + x_2)$ og $q = x_1 \cdot x_2$.
\fo{Sjekk 1} a) $x = 0$ eller $x = -7$ \ b) Ingen, fordi et kvadrat aldri er negativt \
c) Nullproduktregelen gjelder bare når produktet er 0. Riktig: $x^2 + x - 6 = 0$, så $(x+3)(x-2) = 0$ og $x = -3$ eller $x = 2$.

\fk{Del 2 — abc-formelen og diskriminanten}
\fo{2.1} a) $D = -16$, ingen løsning \ b) $D = 16$, to løsninger \ c) $D = 0$, én løsning
\fo{2.2} a) $x = 2$ eller $x = 5$ \ b) $D = 20$, så $x = -1 \pm \sqrt{5}$ \ c) $D = 9$, så $x = -2$ eller $x = -\frac{1}{2}$
\fo{2.3} $x^2 + 4x = 21$ og $x^2 + 4x + 4 = 25$. $(x+2)^2 = 25$, så $x + 2 = \pm 5$. $x = 3$ eller $x = -7$.
\fo{2.4} Metodene under er forslag. Andre valg med begrunnelse er riktige. a) abc-formelen, $D = 12$, så $x = 1 \pm \sqrt{3}$ \ b) abc-formelen, $D = 25$, så $x = \frac{2}{3}$ eller $x = -1$ \
c) abc-formelen, $D = -4$, ingen løsning \ d) kvadratrot: $x^2 = 4$, så $x = 2$ eller $x = -2$
\fo{2.5} Feilen er i overgangen til $\sqrt{-25 + 24}$. Her er $(-5)^2 = 25$, ikke $-25$.
Riktig: $D = 25 + 24 = 49$, $x = \frac{5 \pm 7}{4}$, så $x = 3$ eller $x = -\frac{1}{2}$.
\fo{2.6} Omkretsen $2(x + y) = 34$ gir $x + y = 17$, så den andre siden er $17 - x$.
$x(17 - x) = 60$ gir $x^2 - 17x + 60 = 0$. $D = 49$ og $x = \frac{17 \pm 7}{2}$, så $x = 5$ eller $x = 12$.
Sidene er 5 cm og 12 cm. Kontroll: omkrets $2 \cdot 17 = 34$ og areal $5 \cdot 12 = 60$.
\fo{2.7} $20t - 5t^2 = 15$ gir $t^2 - 4t + 3 = 0$, så $t = 1$ eller $t = 3$.
Ballen er 15 m over bakken på vei opp etter 1 s og på vei ned etter 3 s.
\fo{2.8} $D = 16 - 4k$. a) $k < 4$ \ b) $k = 4$ \ c) $k > 4$ \
d) $D = k^2 - 4 \cdot 1 \cdot 36 = k^2 - 144$. Én løsning krever $D = 0$, så $k^2 = 144$ og $k = 12$ eller $k = -12$. Ingen andre verdier gir $D = 0$.
\fo{2.9} a) Alltid. Er $a$ og $c$ motsatt fortegn, er $ac < 0$, så $-4ac > 0$. Siden $b^2 \ge 0$, er $D = b^2 - 4ac > 0$. To løsninger.
b) Noen ganger. $x^2 + x + 1 = 0$ har $D = -3$, ingen løsning. $x^2 + 5x + 4 = 0$ har $D = 9$ og løsningene $x = -1$ og $x = -4$.
\fo{2.10} Nullpunktene $-1$ og $3$ gir $y = a(x+1)(x-3)$. Ved $x = 0$ er $y = -3a = 6$, så $a = -2$.
$y = -2(x^2 - 2x - 3) = -2x^2 + 4x + 6$. Altså $a = -2$, $b = 4$ og $c = 6$.
\fo{Sjekk 2} a) $D = 9$, så $x = 1$ eller $x = -\frac{1}{2}$ \ b) $D = -3$, så ingen løsning \
c) Grafen tangerer $x$-aksen i ett punkt, så funksjonen har ett nullpunkt.

\fk{Ekstra utfordring}
\fo{U1} a) $u^2 - 5u + 4 = 0$ gir $u = 1$ eller $u = 4$. Da er $x = \pm 1$ eller $x = \pm 2$.
b) $u^2 + u - 6 = (u+3)(u-2) = 0$ gir $u = -3$ eller $u = 2$. Siden $x^2$ ikke kan være $-3$, er $x = \sqrt{2}$ eller $x = -\sqrt{2}$. To løsninger.
\fo{U2} Tillatt: $x \neq 0$, fordi vi ikke kan dele på 0. a) $x^2 + 6 = 5x$ gir $x^2 - 5x + 6 = 0$, så $x = 2$ eller $x = 3$.
b) $x \neq 0$. $2x^2 = 3x - 1$ gir $2x^2 - 3x + 1 = 0$, $D = 1$, så $x = 1$ eller $x = \frac{1}{2}$.
\fo{U3} a) Utskriften er tre linjer: \lstinline[style=laereverk,numbers=none,frame=none]|[-2.0, 4.0]|, \lstinline[style=laereverk,numbers=none,frame=none]|[]| og \lstinline[style=laereverk,numbers=none,frame=none]|[2.0]|. De tre likningene har $D = 36$, $D = -16$ og $D = 0$.
b) Funksjonen deler på $2a = 0$ og gir en feilmelding. Legg inn først i funksjonen
\lstinline[style=laereverk,numbers=none,frame=none]|if a == 0: return [-c / b]|. Da gir \lstinline[style=laereverk,numbers=none,frame=none]|abc(0, 2, -4)| svaret \lstinline[style=laereverk,numbers=none,frame=none]|[2.0]|, forutsatt $b \neq 0$.
\fo{U4} a) $x + 2 = x^2$ gir $x^2 - x - 2 = (x-2)(x+1) = 0$, så $x = 2$ eller $x = -1$.
b) $x = 2$: $\sqrt{4} = 2$, det stemmer. $x = -1$: $\sqrt{1} = 1 \neq -1$, det stemmer ikke.
c) $\sqrt{\phantom{x}}$ er aldri negativ, så $x$ måtte være minst 0. Kvadrering gir ikke det kravet: $(-1)^2 = 1^2$. Fra $A^2 = B^2$ følger bare $A = B$ eller $A = -B$.
\fo{U5} Ett eksempel er $2(x-2)(x+5) = 2x^2 + 6x - 20 = 0$. Flere svar er riktige.
Det finnes uendelig mange. For hvert tall $a \neq 0$ har $a(x-2)(x+5) = 0$ løsningene 2 og $-5$, for hvis produktet er 0 må en av faktorene $(x-2)$ og $(x+5)$ være 0 ($a \neq 0$).
\fo{U6} a) Sum:
\[ x_1 + x_2 = \frac{-b + \sqrt{D}}{2a} + \frac{-b - \sqrt{D}}{2a} = \frac{-2b}{2a} = -\frac{b}{a} \]
Produkt:
\[ x_1 \cdot x_2 = \frac{b^2 - D}{4a^2} = \frac{b^2 - (b^2 - 4ac)}{4a^2} = \frac{c}{a} \]
b) Produktet er $\frac{c}{a} = \frac{-6}{3} = -2$. Da er $3 \cdot x_2 = -2$, så $x_2 = -\frac{2}{3}$.
Kontroll med summen: $3 - \frac{2}{3} = \frac{7}{3} = -\frac{b}{a}$.
